\section{Introduction}

Modern text-to-speech systems increasingly combine a language-model-style
generator with a learned neural speech codec. Qwen3-TTS follows this pattern
and exposes textual VoiceDesign conditioning alongside multilingual synthesis
\cite{hu2026qwen3tts}. Its 12-Hz family predicts a hierarchical set of speech
codes and supports incremental waveform reconstruction. This structure is
well-suited to interactive delivery, but the serving implementation determines
whether model capability translates into bounded memory, low first-audio
latency, and progressive output under concurrent load.

\projectname{} is an inference-system implementation rather than a new model.
It implements the pinned \modelname{} checkpoint as a native Rust/CUDA service
for a specific hardware target: NVIDIA DGX Spark with a GB10 GPU and
\code{sm\_121} machine code. It deliberately narrows the product surface to
textual VoiceDesign. Reference-audio voice cloning, speaker enrollment, the
speech-tokenizer encoder, other Qwen3-TTS variants, and the retired 0.6B model
are out of scope.

The implementation makes four systems contributions:

\begin{enumerate}
  \item a native execution path for the VoiceDesign talker, its 15-step code
        predictor, and the decoder-only neural speech tokenizer;
  \item an event-ordered device-to-device handoff that keeps generated codec
        codes on the GPU until the decoder consumes them;
  \item a fixed-capacity scheduler with request-local CUDA state, bounded PCM
        ownership, backpressure, cancellation, and progressive HTTP framing;
        and
  \item a fail-closed evaluation protocol that compares an immutable native
        candidate with a pinned stock SGLang-Omni comparator using one external
        client and digest-bound raw evidence.
\end{enumerate}

The present source is intentionally conservative. It documents implemented
mechanisms and a complete evaluation design, but does not turn exploratory
measurements into publication claims. Section~\ref{sec:evaluation} defines the
acceptance protocol, while Section~\ref{sec:results} remains evidence-gated.
